\section{Week 2: Coupon Collector, Tail Bounds, Birthday Paradox}\label{sec:week2}

Lectures of September 1 and 3.

\begin{opentable}
    \openrow{def:mclv-0901}{Monte Carlo and Las Vegas algorithms}
    \openrow{thm:coupon}{Coupon collector, $\expectation{X} = n H_n$}
    \openrow{thm:whp}{$3n \log n$ queries access every slot w.h.p.}
    \openrow{thm:lowerbound}{Las Vegas search needs $\bigOm{n}$ (Yao)}
    \openrow{thm:chernoff-0901}{Chernoff bounds (three forms)}
    \openrow{lem:maxload-lower-0901}{Exercise: no slot accessed fewer than $12 \log n$ times}
    \openrow{thm:markov}{Markov's inequality}
    \openrow{thm:chebyshev}{Chebyshev's inequality}
    \openrow{lem:pairwise-0903}{Variance of the number of collisions}
    \openrow{thm:birthday-0903}{Exercise: birthday paradox, $m \geq 2\sqrt{n} + 1$}
    \openrow{thm:collision-lb-0903}{Exercise: collision lower bound $\sfrac{k(k-1)}{2n}$}
\end{opentable}

\paragraph{Which inequality.}
Markov needs $X \geq 0$ and $\expectation{X}$:
$\prob{X \geq b} \leq \sfrac{\expectation{X}}{b}$. Chebyshev needs the variance:
$\prob{\abs{X - \expectation{X}} \geq b} \leq \sfrac{\variance{X}}{b^2}$. Chernoff needs a sum of
independent $0/1$ variables with mean $\mu$:
\begin{gather*}\prob{X \geq (1+\delta)\mu} \leq e^{-\mu\delta^2/3},
    \qquad
    \prob{X < (1-\delta)\mu} \leq e^{-\mu\delta^2/2}
    \quad (0 < \delta < 1), \\
    \prob{X \geq R} \leq 2^{-R} \quad (R \geq 6\mu).\end{gather*}
For a statement about all $n$ bins, slots or coupons, bound one of them and multiply by $n$ (union
bound).

\subsection{Tail Bounds}

\begin{exercisebox}
\begin{pproblem}[The same event, three bounds]\label{pp:three}
    A fair coin is tossed $n$ times and $X$ is the number of heads. Bound
    $\prob{X \geq \sfrac{3n}{4}}$ with each inequality. Evaluate for $n = 100$.
\end{pproblem}
\tcblower
\begin{wsolution}
    \uses{\scr{thm:markov}; \scr{thm:chebyshev}; \scr{thm:chernoff-0901}}
    $\expectation{X} = \sfrac{n}{2}$ and $\variance{X} = n \cdot \frac{1}{2} \cdot \frac{1}{2} = \sfrac{n}{4}$.
    \begin{enumerate}
        \item Markov: $\prob{X \geq \sfrac{3n}{4}} \leq \frac{n/2}{3n/4} = \frac{2}{3}$.
        \item Chebyshev: $X \geq \sfrac{3n}{4}$ implies $\abs{X - \sfrac{n}{2}} \geq \sfrac{n}{4}$, so
              the probability is at most $\frac{n/4}{(n/4)^2} = \frac{4}{n}$.
        \item Chernoff: $\sfrac{3n}{4} = \left(1 + \frac{1}{2}\right)\mu$ with $\mu = \sfrac{n}{2}$ and
              $\delta = \frac{1}{2}$, so the probability is at most
              $e^{-(n/2)(1/4)/3} = e^{-n/24}$.
    \end{enumerate}
    For $n = 100$: $0.667$, $0.04$ and $e^{-4.17} \approx 0.016$. Markov does not improve with $n$,
    Chebyshev improves like $\sfrac{1}{n}$, and Chernoff improves exponentially.
\end{wsolution}
\end{exercisebox}

\begin{exercisebox}
\begin{pproblem}[$n$ balls into $n$ bins]\label{pp:balls}
    Throw $n$ balls independently and uniformly into $n$ bins, $n \geq 4$. Show that with probability
    at least $1 - \sfrac{1}{n^2}$ no bin receives $3 \log_2 n$ balls or more.
\end{pproblem}
\tcblower
\begin{wsolution}
    \uses{\scr{thm:chernoff-0901}}
    \leavevmode
    \begin{enumerate}
        \item Fix a bin and let $X$ be its number of balls. $X$ is a sum of $n$ independent
              indicators, each $1$ with probability $\sfrac{1}{n}$, so $\mu = 1$.
        \item Take $R = 3 \log_2 n$. For $n \geq 4$ we have $R \geq 6 = 6\mu$, so the third Chernoff
              form applies.
        \item $\prob{X \geq 3 \log_2 n} \leq 2^{-3 \log_2 n} = \sfrac{1}{n^3}$.
        \item Union bound over the $n$ bins:
              $\prob{\text{some bin has at least } 3 \log_2 n \text{ balls}} \leq n \cdot \frac{1}{n^3} = \frac{1}{n^2}$. \qedhere
    \end{enumerate}
\end{wsolution}
\end{exercisebox}


\begin{exercisebox}
\begin{pproblem}[Max load, upper tail, with the scribe's forms]\label{wp:maxload}
    Random sampling makes $3n \ln n$ queries into an array of $n \geq 2$ slots. Show that with
    probability at least $1 - n^{-11}$ no slot is accessed $18 \ln n$ times or more.
\end{pproblem}
\tcblower
\begin{wsolution}
    \uses{\scr{thm:chernoff-0901}; \scr{lem:maxload-lower-0901} for the lower tail}
    \leavevmode
    \begin{enumerate}
        \item Fix a slot and let $Y$ be the number of queries that access it. $Y$ is a sum of
              $3n \ln n$ independent indicators, each $1$ with probability $\sfrac{1}{n}$, so
              $\mu = 3 \ln n$.
        \item $R = 18 \ln n = 6\mu$, so the third form applies:
              $\prob{Y \geq 18 \ln n} \leq 2^{-18 \ln n} = n^{-18 \ln 2} \leq n^{-12}$, since
              $18 \ln 2 \approx 12.48$.
        \item Union bound over the $n$ slots: the failure probability is at most
              $n \cdot n^{-12} = n^{-11}$. \qedhere
    \end{enumerate}
    \emph{Beyond the scribes.} The professor's notes state the sharper cap $9 \ln n$. That needs the
    general Chernoff form
    $\prob{X \geq (1+\delta)\mu} \leq \left(\sfrac{e^\delta}{(1+\delta)^{1+\delta}}\right)^\mu$ with
    $\delta = 2$, which the scribe does not state (proof in \texttt{practiceproblems.pdf},
    Sections 10--11).
\end{wsolution}
\end{exercisebox}

\begin{exercisebox}
\begin{pproblem}[Coupon collector with a different number of trials]\label{pp:coupon}
    There are $n$ coupons and each trial gives a uniformly random one. Bound the probability that
    some coupon is missing after
    \begin{enumerate}[label=(\alph*)]
        \item $2n \ln n$ trials,
        \item $n \ln n + cn$ trials.
        \item What happens after only $n$ trials?
    \end{enumerate}
\end{pproblem}
\tcblower
\begin{wsolution}
    \uses{\scr{thm:whp}; \scr{thm:linearity-0827}}
    A fixed coupon is missed by one trial with probability $1 - \sfrac{1}{n}$, so it is missing
    after $m$ trials with probability $\left(1 - \frac{1}{n}\right)^m \leq e^{-m/n}$.

    (a) $e^{-2 \ln n} = \sfrac{1}{n^2}$ per coupon, so at most $n \cdot \frac{1}{n^2} = \frac{1}{n}$
    by the union bound.

    (b) $e^{-\ln n - c} = \sfrac{e^{-c}}{n}$ per coupon, so at most $e^{-c}$ in total. For $c = 3$
    this is about $0.05$.

    (c) The expected number of missing coupons is
    $n\left(1 - \frac{1}{n}\right)^n \approx \sfrac{n}{e}$, by linearity of expectation. About
    $37\%$ of the coupons are still missing.
\end{wsolution}
\end{exercisebox}

\begin{exercisebox}
\begin{pproblem}[Boosting a Monte Carlo algorithm by majority vote]\label{pp:majority}
    An algorithm gives the correct yes/no answer with probability $\sfrac{2}{3}$. Run it $t$ times
    independently ($t$ odd) and output the majority answer. How large must $t$ be for an error of at
    most $\sfrac{1}{n}$?
\end{pproblem}
\tcblower
\begin{wsolution}
    \uses{\scr{def:mclv-0901}; \scr{thm:chernoff-0901}}
    \leavevmode
    \begin{enumerate}
        \item Let $X$ be the number of correct runs: a sum of $t$ independent indicators with
              $\mu = \sfrac{2t}{3}$.
        \item The majority is wrong only if $X < \sfrac{t}{2}$.
        \item Write $\sfrac{t}{2} = (1 - \delta)\mu$: $(1 - \delta) \cdot \frac{2t}{3} = \frac{t}{2}$
              gives $\delta = \frac{1}{4}$.
        \item Lower tail:
              $\prob{X < \sfrac{t}{2}} \leq e^{-\mu\delta^2/2} = e^{-(2t/3)(1/16)/2} = e^{-t/48}$.
        \item $e^{-t/48} \leq \sfrac{1}{n}$ when $t \geq 48 \ln n$.
    \end{enumerate}
    So $\bigO{\log n}$ repetitions bring the error from $\sfrac{1}{3}$ down to $\sfrac{1}{n}$.
\end{wsolution}
\end{exercisebox}

\begin{exercisebox}
\begin{pproblem}[Empty bins]\label{pp:empty}
    Throw $m$ balls independently and uniformly into $n$ bins.
    \begin{enumerate}[label=(\alph*)]
        \item Find the expected number of empty bins.
        \item Evaluate it for $m = n$, $m = n \ln n$ and $m = 2n \ln n$.
        \item For $m = 2n \ln n$, bound the probability that some bin is empty.
    \end{enumerate}
\end{pproblem}
\tcblower
\begin{wsolution}
    \uses{\scr{thm:linearity-0827}; \scr{thm:markov}}
    Let $Y_i = 1$ if bin $i$ is empty and $0$ otherwise, and let $Y = Y_1 + \dots + Y_n$ be the
    number of empty bins. One ball misses bin $i$ with probability $1 - \sfrac{1}{n}$, and the balls
    are independent, so
    $\expectation{Y_i} = \prob{Y_i = 1} = \left(1 - \frac{1}{n}\right)^m$.

    (a) By linearity of expectation,
    $\expectation{Y} = n\left(1 - \frac{1}{n}\right)^m \leq n e^{-m/n}$, using $1 - x \leq e^{-x}$.

    (b) For $m = n$: about $\sfrac{n}{e}$, so $37\%$ of the bins are empty (for $n = 100$ the exact
    value is $36.6$). For $m = n \ln n$: at most $n e^{-\ln n} = 1$. For $m = 2n \ln n$: at most
    $n \cdot n^{-2} = \sfrac{1}{n}$.

    (c) $Y \geq 0$, so Markov gives $\prob{Y \geq 1} \leq \expectation{Y} \leq \sfrac{1}{n}$.

    The $Y_i$ are not independent (if every other bin is empty, bin $i$ is not), and linearity does
    not need them to be. Part (c) is the coupon collector bound again: Markov on the number of empty
    bins gives the same number as the union bound over the bins.
\end{wsolution}
\end{exercisebox}

\begin{exercisebox}
\begin{pproblem}[Las Vegas to Monte Carlo]\label{pp:lv-mc}
    A Las Vegas algorithm $\mathcal{A}$ is always correct and makes $T$ queries in expectation on
    every input.
    \begin{enumerate}[label=(\alph*)]
        \item Build an algorithm that never makes more than $2T$ queries and fails with probability
              at most $\frac{1}{2}$.
        \item Get the failure probability down to $\sfrac{1}{n}$.
        \item Check (a) on random sampling for search.
    \end{enumerate}
\end{pproblem}
\tcblower
\begin{wsolution}
    \uses{\scr{def:mclv-0901}; \scr{thm:markov}}
    (a) Run $\mathcal{A}$ and stop it after $2T$ queries. If it has finished, output its answer,
    which is correct. If not, output ``fail''. Let $X$ be the number of queries that $\mathcal{A}$
    would make. $X \geq 0$ and $\expectation{X} \leq T$, so by Markov
    $\prob{X \geq 2T} \leq \frac{\expectation{X}}{2T} \leq \frac{1}{2}$. The new algorithm fails only
    in that case. Stopping after $cT$ queries gives failure at most $\sfrac{1}{c}$ in the same way.

    (b) Markov alone would need $c = n$, that is $nT$ queries. Instead repeat the $2T$-query version
    $\ceil{\log_2 n}$ times with fresh randomness and stop at the first run that finishes. All runs
    fail with probability at most $\left(\frac{1}{2}\right)^{\log_2 n} = \frac{1}{n}$. The total is
    at most $2T \ceil{\log_2 n}$ queries.

    (c) Random sampling with replacement finds $x$ with probability $\sfrac{1}{n}$ per query, so the
    number of queries is geometric with mean $T = n$. Markov says that $2n$ queries fail with
    probability at most $\frac{1}{2}$. The exact value is
    $\left(1 - \frac{1}{n}\right)^{2n} \leq e^{-2} \approx 0.135$.

    Markov is weak here, but it used nothing about $\mathcal{A}$ except its expected cost. Part (b)
    works because a Las Vegas run tells us whether it finished, so failures are detected and never
    output as answers.
\end{wsolution}
\end{exercisebox}


\subsection{Lower Bound for Randomized Search}

\begin{exercisebox}
\begin{pproblem}[Randomized search needs about $n/2$ expected queries]\label{pp:yao-lv}
    Use Yao's principle to show that a Las Vegas algorithm for search makes at least
    $\sfrac{(n-1)}{2}$ expected queries on some input.
\end{pproblem}
\tcblower
\begin{wsolution}
    \uses{\scr{thm:lowerbound}; \scr{thm:detsearch-0825}}
    \leavevmode
    \begin{enumerate}
        \item Input distribution $\mu$: $x$ is in a uniformly random slot, each with probability
              $\sfrac{1}{n}$.
        \item Take a deterministic algorithm that is always correct. As long as every answer is
              ``no'', its queries follow one fixed order $q_1, q_2, \dots$
        \item It cannot stop while two or more slots are unqueried and $x$ has not been seen (the
              adversary argument). So if $x$ is in slot $q_i$ with $i \leq n - 1$, it makes at least
              $i$ queries, and if $x$ is in the last slot it makes $n - 1$.
        \item The expected number of queries under $\mu$ is at least
              \[\frac{1}{n} \left(\sum_{i=1}^{n-1} i + (n-1)\right) = \frac{(n-1)(n+2)}{2n} \geq \frac{n-1}{2}.\]
        \item By Yao's principle the expected cost of a randomized algorithm on its worst input is
              at least the expected cost of the best deterministic algorithm under $\mu$. So it is
              at least $\sfrac{(n-1)}{2} = \bigOm{n}$. \qedhere
    \end{enumerate}
\end{wsolution}
\end{exercisebox}


\subsection{Birthday Paradox and Collisions}

There are $m$ people (queries) and $n$ days (values), and $X$ is the number of colliding pairs:
\begin{gather*}
    \expectation{X} = \frac{m(m-1)}{2n},
    \qquad
    \prob{X = 0} \leq \frac{\variance{X}}{\expectation{X}^2} \leq \frac{1}{\expectation{X}}, \\
    \prob{\text{no collision}} = \prod_{i=0}^{m-1} \left(1 - \frac{i}{n}\right) \leq e^{-m(m-1)/(2n)}.
\end{gather*}

\begin{exercisebox}
\begin{pproblem}[The numbers for $n = 365$]\label{pp:bday365}
    \leavevmode
    \begin{enumerate}[label=(\alph*)]
        \item How many people make the expected number of collisions at least $1$?
        \item How many make the probability of a collision at least $\frac{1}{2}$, by the chain-rule
              bound?
        \item How many does the Chebyshev proof need to guarantee probability $\frac{1}{2}$?
    \end{enumerate}
\end{pproblem}
\tcblower
\begin{wsolution}
    \uses{\scr{thm:birthday-0903}}
    (a) We need $m(m-1) \geq 2n = 730$. Since $27 \cdot 26 = 702$ and $28 \cdot 27 = 756$, the
    answer is $m = 28$, with $\expectation{X} = \sfrac{756}{730} \approx 1.04$.

    (b) We need $e^{-m(m-1)/730} \leq \frac{1}{2}$, that is $m(m-1) \geq 730 \ln 2 \approx 505.997$.
    Since $22 \cdot 21 = 462$ and $23 \cdot 22 = 506$, the answer is $m = 23$. The exact product for
    $m = 23$ is $0.4927$, so a collision has probability $0.507$.

    (c) Chebyshev gives $\prob{X = 0} \leq \sfrac{1}{\expectation{X}}$, so we need
    $\expectation{X} \geq 2$, that is $m(m-1) \geq 1460$. Since $38 \cdot 37 = 1406$ and
    $39 \cdot 38 = 1482$, the answer is $m = 39$. The condition $m \geq 2\sqrt{n} + 1 \approx 39.2$
    of the exercise gives $m = 40$.

    The three answers differ ($28$, $23$, $39$) but all are about $\sqrt{n}$. Part (a) alone proves
    nothing about the probability: $\expectation{X} \geq 1$ does not imply that
    $\prob{X \geq 1}$ is large. That step needs more than the expectation: either the direct
    product of part (b), or the variance, which is what the exercise uses with Chebyshev.
\end{wsolution}
\end{exercisebox}

\begin{exercisebox}
\begin{pproblem}[Why Chebyshev and not Chernoff?]\label{pp:why-cheb}
    $X$ is a sum of indicators $X_{ij}$. Why does the birthday proof not use the Chernoff bound?
    Why is Markov not enough?
\end{pproblem}
\tcblower
\begin{wsolution}
    \uses{\scr{lem:pairwise-0903}; \scr{thm:birthday-0903}}
    \leavevmode
    \begin{enumerate}
        \item Chernoff needs mutually independent indicators. The $X_{ij}$ are not: if persons $1$
              and $2$ share a birthday and persons $2$ and $3$ share a birthday, then persons $1$
              and $3$ do too. So $\prob{X_{12} = X_{23} = X_{13} = 1} = \sfrac{1}{n^2}$, not
              $\sfrac{1}{n^3}$.
        \item They are pairwise independent: for two different pairs,
              $\prob{X_{ij} = X_{kl} = 1} = \sfrac{1}{n^2} = \prob{X_{ij} = 1} \cdot \prob{X_{kl} = 1}$,
              whether the pairs share a person or not.
        \item Pairwise independence is exactly what the variance needs. The cross terms
              $\expectation{X_{ij} X_{kl}} - \expectation{X_{ij}}\expectation{X_{kl}}$ are all zero,
              so $\variance{X}$ is the sum of the $\variance{X_{ij}}$, which is at most
              $\expectation{X}$. Chebyshev then applies.
        \item Markov gives $\prob{X \geq 1} \leq \expectation{X}$. That is an upper bound on the
              probability of a collision, and we want a lower bound.
    \end{enumerate}
    The rule: Markov needs only $X \geq 0$, Chebyshev needs the variance (pairwise independence is
    enough), and Chernoff needs mutual independence.
\end{wsolution}
\end{exercisebox}

\begin{exercisebox}
\begin{pproblem}[The collision lower bound with numbers]\label{pp:collision-num}
    Let $f$ be a random function into $n = 1{,}000{,}000$ values.
    \begin{enumerate}[label=(\alph*)]
        \item Bound the probability that $500$ queries find a collision.
        \item In general, bound it for $k = \sfrac{\sqrt{n}}{2}$ queries.
        \item How many queries are needed for probability $\frac{1}{2}$, and how many are enough?
    \end{enumerate}
\end{pproblem}
\tcblower
\begin{wsolution}
    \uses{\scr{thm:collision-lb-0903}; \scr{thm:birthday-0903}}
    By the Sep 3 exercise, $k$ queries find a collision with probability at most
    $\sfrac{k(k-1)}{(2n)}$. The reason is that query number $s$ hits one of the $s - 1$ earlier
    answers with probability at most $\sfrac{(s-1)}{n}$, and the union bound adds these up.

    (a) $\frac{500 \cdot 499}{2{,}000{,}000} \approx 0.125$.

    (b) $\frac{k(k-1)}{2n} < \frac{k^2}{2n} = \frac{n/4}{2n} = \frac{1}{8}$, for every $n$.

    (c) Needed: $\frac{k(k-1)}{2n} \geq \frac{1}{2}$ means $k(k-1) \geq 1{,}000{,}000$. Since
    $1000 \cdot 999 = 999{,}000$ is too small, $k \geq 1001$. Enough: by the birthday exercise,
    $m \geq 2\sqrt{n} + 1 = 2001$ random queries collide with probability at least $\frac{1}{2}$.

    So the answer lies between $1001$ and $2001$ queries: $\bigTh{\sqrt{n}}$. These are two
    separate proofs. The upper bound analyses one algorithm (random queries, Chebyshev). The lower
    bound holds for every algorithm, adaptive or randomized (union bound).
\end{wsolution}
\end{exercisebox}

